\begin{definition} \label{definition:orientation_of_a_topological_manifold_with_or_without_boundary_using_singular_homology}
Let $M$ be a \CrefAndHyperref{definition:topological_manifold}{topological $n$-manifold (with or without boundary)}.

An \hldef{orientation of $M$} is a function $x \mapsto \mu_x$ assigning a \CrefAndHyperref{definition:local_orientation_of_a_topological_manifold_with_or_without_boundary_at_an_interior_point_using_local_singular_homology}{local orientation} to each point $x \in \operatorname{int} M$ \CrefIfExists{definition:interior_and_boundary_of_topological_space} such that for every $x \in \operatorname{int} M$, there exists an open neighborhood $U \subset \operatorname{int} M$ and a class $\alpha \in H_n(M, M \setminus U; \mathbb{Z})$ such that for all $y \in U$, the map induced by the inclusion $j_y: (M, M \setminus U) \to (M, M \setminus \{y\})$ sends $\alpha$ to $\mu_y$.

If $M$ is a topological manifold with or without boundary for which an orientation exists, then $M$ is said to be \hldef{orientable}.

A topological manifold $M$ with or without boundary equipped with an orientation is said to be an \hldef{oriented manifold}. 
\end{definition}